% DER_XIV: Grothendieck's text with the commentary set in, then the overview, inside a group that applies the chunk's own definitions
\dgrchapter{XIV}{Chapter XIV}{Carrés h-cartésiens et h-cocartésiens}
\dgrpart{0}{XIV}{Manuscript pages 1--8}{XIV\enspace\textperiodcentered\enspace Chapter XIV\enspace\textperiodcentered\enspace manuscript pages 1--8}{Chapter XIV \textperiodcentered\ \textup{XIV}}
\dgrchunkbegin
% DER_XIV: the definitions of the chunk's preamble, applied inside its group
\let\printedpage\relax\newcommand{\printedpage}[1]{\if\relax\detokenize{#1}\relax\else\begingroup\reversemarginpar\marginnote{\footnotesize p.~#1}\endgroup\fi}
\let\msp\relax\newcommand{\msp}[1]{\marginnote{\footnotesize [#1]}}

% footnote marks as printed: "(1)" in the text, "1" before the footnote text
\makeatletter
\providecommand\@makefnmark{}\renewcommand\@makefnmark{\hbox{\textsuperscript{\normalfont(\@thefnmark)}}}
\providecommand\@makefntext{}\renewcommand\@makefntext[1]{\parindent 1em\noindent\hb@xt@1.8em{\hss\textsuperscript{\normalfont\@thefnmark}}#1}
\makeatother
\let\uHom\relax\newcommand{\uHom}{\underline{\mathrm{Hom}}}
\let\bDelta\relax\newcommand{\bDelta}{\mathbf{\Delta}}
\let\ho\relax\newcommand{\ho}{\mathrm{ho}}
\let\cM\relax\newcommand{\cM}{\mathcal{M}}

% ---- from Der-ChXIV-07-10-10_p001-005.tex
% ===================== PDF page 1 =====================
\printedpage{}% no page number printed on this page
% head: (none)

\begin{center}
\vspace*{6cm}

{\LARGE \textsf{LES DÉRIVATEURS}}

\vspace{1cm}

{\large \textbf{ALEXANDRE GROTHENDIECK}}

\vspace{2cm}

{\large \textbf{Chapitre XIV}}

\vspace{0.3cm}

{\large \textbf{Carrés h-cartésiens et h-cocartésiens}}
\end{center}

\clearpage

% ===================== PDF page 2 =====================
\printedpage{}% no page number printed on this page
% head: (none)

\vspace*{6cm}

\noindent Ce texte a été déchiffré et transcrit en \LaTeXe{} par M. Künzer. Il a été édité par M. Künzer, J. Malgoire, et G. Maltsiniotis. La transcription est aussi fidèle que possible au manuscrit. Pour les quelques corrections évidentes, ou rares commentaires des éditeurs, ainsi que pour la numérotation originale des pages du manuscrit, les caractères de machine à écrire \texttt{[typewriter]} entre crochets sont utilisés. Un point d'interrogation entre crochets signifie que l'on n'est pas sûr du mot qui précède.

\noindent Ce chapitre est inachevé. Dans la première section, Grothendieck introduit une notion de “catégorie quadrangulée” à gauche (resp. à droite). Le texte s'interrompt brutalement au début de la deuxième section, consacrée aux exemples, quand Grothendieck réalise, sans doute, que l'exemple principal qu'il a en vue (la catégorie homotopique d'une catégorie de modèles munie des carrés homotopiquement cartésiens (resp. cocartésiens)) ne satisfait \textit{pas} à l'axiomatique des catégories quadrangulées. Cela est brièvement expliqué dans une postface des éditeurs.

\noindent Cette édition est provisoire. Les remarques, commentaires et corrections sont bienvenus. Envoyer un message à :

\noindent \textbf{[e-mail address removed]}

\hfill G. Maltsiniotis\hspace*{2cm}

\clearpage

% ===================== PDF page 3 =====================
\printedpage{1}
% head: A. Grothendieck : Dérivateurs, Ch. XIV, Édition des Universités de Montpellier II et Paris VII
% foot: Version 7 octobre 2010

\msp{page 1}

\section*{1\quad Axiomes de structures quadrangulées (une tentative avortée)}

Soit $(\mathcal{M},W)$ une catégorie de modèles admettant une théorie de la fibre homotopique. Alors on a dans $\mathrm{Ho}_W=\mathcal{M}W^{-1}$ une notion de \textit{carré $W$-cartésien} (ou $W$-homotopiquement cartésien), par quoi on entend un diagramme carré
\[
\begin{tikzcd}
X_1 \arrow[d] & X_3 \arrow[l] \arrow[d] \\
X_0 & X_2 \arrow[l]
\end{tikzcd}
\]
dans $\mathrm{Ho}_W$ provenant, à isomorphisme près, d'un carré $W$-cartésien dans $\mathcal{M}$ – donc isomorphe, dans $\underline{\mathrm{Hom}}(\boldsymbol{\Delta}_1\times\boldsymbol{\Delta}_1,\mathrm{Ho}_W)$, a un carré \textit{cartésien} \texttt{[dans $\mathcal{M}$]} % [sic] a un carré -> à un carré, preposition printed without its grave accent
\[
\begin{tikzcd}
Y_1 \arrow[d, "p_1"'] & Y_3\ (\simeq Y_1\times_{Y_0}Y_2) \arrow[l] \arrow[d] & \\
Y_0 & Y_2 \arrow[l, "p_2"] & ,
\end{tikzcd}
\]
avec $p_1,\ p_2\in\mathrm{Fib}_W$. (En fait, il suffit de $p_1$ ou $p_2$ dans $\mathrm{Fib}_W$ pour assurer le caractère $W$-cartésien du carré cartésien envisagé.) Je voudrais axiomatiser cette situation, histoire surtout de faire la liste des propriétés serviables essentielles de cette notion, et de la notion duale de carré $W$-cocartésien (ou $W$-homotopiquement cocartésien), dans le cas où $(\mathcal{M},W)$ admet une théorie de la cofibre homotopique. L'espèce de structure, savoir celle des “catégories quadrangulées” (à gauche, à droite, bilatères) est une variante, dans un cas \msp{page 2} non additif, de la notion de catégorie triangulée dégagée par \textsc{Verdier}. Comme celle-ci, cette notion est en fait trop faible – ce n'est \textit{pas} “la bonne” – mais elle a un intérêt pratique provisoire. La “bonne” notion (s'il en est une à part celle de dérivateur) devrait travailler avec une notion plus fine de carré et de carré $W$-cartésien, en travaillant systématiquement dans $\mathrm{HO}_W(\boldsymbol{\Delta}_1\times\boldsymbol{\Delta}_1)=\mathcal{M}(\boldsymbol{\Delta}_1\times\boldsymbol{\Delta}_1)W(\boldsymbol{\Delta}_1\times\boldsymbol{\Delta}_1)^{-1}$ (où vivent les “vrais” carrés $W$-cartésiens !), et catégories similaires, telles que $\mathrm{HO}_W(\Psi)$, $\mathrm{HO}_W(\Phi)$. Mais les inclure dans une axiomatique ne faisant pas mention d'une catégorie de modèles $(\mathcal{M},W)$, revient manifestement à mettre le doigt dans l'engrenage de la notion de dérivateur, et ce n'est pas mon propos pour l'instant.

\dgrblock{dgrb741}{What a square must remember}
\begin{dcomment}{exposition}{p1}{What a square must remember}
Grothendieck wants a usable calculus of homotopy pullbacks on the category of homotopy types. The attraction is clear: four objects and four arrows are much easier to manipulate than a model category. The difficulty is already identified in his next paragraph. The comparison
\[
 \operatorname{dia}_{\square}:
 \mathcal M^{\square}(W^{\square})^{-1}
 \longrightarrow (\mathcal MW^{-1})^{\square}
\]
forgets part of a diagram. Here $\square$ has arrows $3\to1\to0$ and $3\to2\to0$, with the two composites equal; $\Phi\subset\square$ omits $3$. We use covariant diagrams, following the arrows actually drawn. An object on the left is a \emph{coherent square}; on the right it is an ordinary commutative square in the localized category. Calling an object coherent does not mean that its chosen representative has stopped commuting strictly. It means that the whole diagram was localized together.

Write $\operatorname{Cart}_{\mathrm{im}}$ for the replete image, on the right, of the homotopy-cartesian objects on the left. Membership says that \emph{some} coherent realization is cartesian. Proposition~1 will assert, incorrectly, that \emph{every} strict realization is then cartesian. An isomorphism of underlying squares need not lift to an isomorphism of coherent squares. Likewise, four arrows forming a cube after localization need not form a morphism of the chosen coherent squares. These are the two losses that the postface makes explicit.

The opening fibration formula belongs to the assumed theory of homotopy fibres. For an arbitrary Quillen model category, a safe version is that $Y_1\times_{Y_0}Y_2$ computes the homotopy pullback when $Y_1\to Y_0$ is a fibration and $Y_0,Y_2$ are fibrant; right properness permits the familiar version without these fibrancy assumptions \citep[Theorem~5.2.5 and the mixed Reedy argument of Lemma~5.2.6, dualized]{Hovey1999ModelCategories}. Thus the opening paragraph should not be detached from its hypotheses or used to diagnose a false formula in the transcription. The editors state the precise Quillen version below.

The commentary on Chapter~I, manuscript p.~7 (``Coherent squares and what the axioms remember''), and on Chapter~IV, manuscript p.~75 (``Constructing the loop square without choosing homotopies''), supplies the general background. The new issue here is reflection of a universal property through $\operatorname{dia}_{\square}$, and the chapter itself supplies a counterexample.
\end{dcomment}


\noindent\textbf{Définition 1.} J'appelle \textit{“catégorie à carrés h-cartésiens”,} ou \textit{“catégorie quadrangulée à gauche”,} une catégorie $\mathcal{H}$, munie d'un sous-ensemble Cart de l'ensemble $\mathrm{Hom}(\boldsymbol{\Delta}_1\times\boldsymbol{\Delta}_1,\mathcal{H})$
\printedpage{2}
% head: A. Grothendieck : Dérivateurs, Ch. XIV, Édition des Universités de Montpellier II et Paris VII
% foot: Version 7 octobre 2010
des diagrammes carrés \texttt{[par carré on entend carré commutatif]}
\[
\begin{tikzcd}
X_1 \arrow[d] & X_3 \arrow[l] \arrow[d] \\
X_0 & X_2 \arrow[l]
\end{tikzcd}
\]
dans $\mathcal{H}$, carrés appelés \textit{“h-cartésiens”} (h comme “homotopiquement”), cette \msp{page 3} donnée étant soumise aux axiomes qui suivent, Cart (1) à Cart (\ldots).

\dgrblock{dgrb742}{The earlier reason for keeping diagram categories}
\begin{dcomment}{context}{p2}{The earlier reason for keeping diagram categories}
The retreat to diagram categories on manuscript page~2 is part of a long-standing programme. In \emph{Pursuing Stacks}, \S69, Grothendieck relates it to the difficulty of formulating trace additivity and determinant multiplicativity using only a derived category. His proposed object retains the derived categories of diagram categories, including filtered complexes of varying lengths \citep[pp.~1--2 of the edited excerpt]{Grothendieck1983}. The letter to Thomason of 2~April~1991 returns to the mapping-cone problem and the need for this whole family \citep[\S2, pp.~5--6]{Grothendieck1991Letter}. These passages explain why an apparently modest square calculus leads back to derivators.

Heller's \emph{Homotopy Theories} already organizes homotopy Kan extensions on diagram categories \citep[Introduction, pp.~1--2]{Heller1988}. This is a historical parallel, not a warrant for all the proposed square rules. The point of the present fragment is precisely to test how much of that structure can survive on the single category $\mathcal H$.

Two different judgments in the source must therefore be distinguished. The ordinary square is too impoverished to carry the desired coherent universal property; nevertheless, some of the axioms imposed on such squares are too strong for their intended examples. There is no contradiction between these diagnoses. Nor should the editors' suggestion that Grothendieck stopped when he noticed the failure be converted into a documented account of his decision: the editors' postface proves the mathematical obstruction, not that biographical explanation.
\end{dcomment}


\noindent Comme je ne suis pas sûr de connaître d'avance toutes les propriétés essentielles, le nombre de ces axiomes reste pour le moment ouvert.

\noindent\textbf{Cart (1) (Symétrie).} Si un carré
\[
\begin{tikzcd}
X_1 \arrow[d] & X_3 \arrow[l] \arrow[d] \\
X_0 & X_2 \arrow[l]
\end{tikzcd}
\]
est h-cartésien, de même le carré symétrique (obtenu en échangeant le rôle de $X_1$ et de $X_2$).

\noindent\textbf{Cart (2) (Carrés triviaux).} Soit
\[
\begin{tikzcd}
X_1 \arrow[d, "p_1"'] & X_3 \arrow[l, "q_2"'] \arrow[d, "q_1"] \\
X_0 & X_2 \arrow[l, "p_2"]
\end{tikzcd}
\]
un carré dans $\mathcal{H}$, avec $p_1$ un isomorphisme. Alors le carré est h-cartésien si et seulement si $q_1$ est un isomorphisme.

\noindent\textbf{Cart (3) (Carrés composés).} Soit un diagramme commutatif
\[
\begin{tikzcd}
X'_0 \arrow[d] \arrow[dr, phantom, "\mathrm{I}"] & X'_1 \arrow[l] \arrow[d] \arrow[dr, phantom, "\mathrm{II}"] & X'_2 \arrow[l] \arrow[d] & \\
X_0 & X_1 \arrow[l] & X_2 \arrow[l] & ,
\end{tikzcd}
\]
d'où les carrés I, II et III :
\[
\begin{tikzcd}
X'_0 \arrow[d] & X'_2 \arrow[l] \arrow[d] \\
X_0 & X_2 \arrow[l]
\end{tikzcd}
\]
\printedpage{3}
% head: A. Grothendieck : Dérivateurs, Ch. XIV, Édition des Universités de Montpellier II et Paris VII
% foot: Version 7 octobre 2010
(qu'on appelle \textit{carré composé} (“horizontal”) des carrés I et II). Supposons I h-cartésien. Alors II est h-cartésien si et seulement si III l'est \texttt{[voir postface des éditeurs]}.

\noindent\textbf{Corollaire.} \textit{Soit un diagramme commutatif cubique, qu'on peut considérer comme homomorphisme d'un carré} (\textit{face supérieure du cube}) \textit{dans un autre} (\textit{face inférieure}) :
\[
\begin{tikzcd}
 & {} \arrow[dl] \arrow[dd] & & {} \arrow[ll] \arrow[dl] \arrow[dd] \\
{} \arrow[dd] & & {} \arrow[ll, crossing over] & \\
 & {} \arrow[dl] & & {} \arrow[ll] \arrow[dl] \\
{} & & {} \arrow[ll] \arrow[from=uu, crossing over] & .
\end{tikzcd}
\]
\textit{Si toutes les faces du cube}\footnote{En fait, il suffit de la face inférieure, et des deux faces parallèles latérales (p. ex. celle de devant et celle de derrière).} \textit{sauf la face supérieure sont} \msp{page 4} \textit{des carrés h-cartésiens (savoir la face inférieure et les quatre faces latérales), il en est de même de la face supérieure.}\footnote{\textbf{Corollaire} (de Cart (3) et Cart (2)) : \textit{tout carré isomorphe à un carré h-cartésien est h-cartésien.}}

\dgrblock{dgrb743}{Pasting, cancellation, and the cube}
\begin{dcomment}{exposition}{p3}{Pasting, cancellation, and the cube}
The first three axioms try to isolate the part of pullback calculus that can be used without constructing a pullback. Cart~(2) has exactly the usual orientation: if the left vertical arrow is invertible, a square is distinguished precisely when the right vertical arrow is invertible. Cart~(1) supplies the transposed assertion. Cart~(3) then includes both composition and cancellation, with the left square held cartesian. Its failure for $\operatorname{Cart}_{\mathrm{im}}$ will be a failure of cancellation; the postface should not be read as disproving each implication separately.

We can fill in the cube argument entirely formally. Regard the cube as two horizontal squares $S_X,S_Y$ and name the opposite lateral faces
\[
 B=\left(\begin{matrix}X_1&\leftarrow&X_3\\\downarrow&&\downarrow\\Y_1&\leftarrow&Y_3\end{matrix}\right),
 \qquad
 F=\left(\begin{matrix}X_0&\leftarrow&X_2\\\downarrow&&\downarrow\\Y_0&\leftarrow&Y_2\end{matrix}\right).
\]
If $B$ and $S_Y$ are distinguished, their vertical composite $R$ is distinguished. The same rectangle is obtained by composing $S_X$ with $F$. If $F$ is distinguished, vertical cancellation implies that $S_X$ is distinguished. Vertical pasting is horizontal pasting after Cart~(1). This proves exactly the stronger three-face version in the first footnote; the other two lateral faces are unnecessary. An isomorphism of squares makes $B,F$ distinguished by Cart~(2), so the second footnote follows as well.

There is also an economy in the later list: Cart~(5)(b) is already a consequence of Cart~(1)--(3). For a cube between distinguished $S_X,S_Y$, suppose $p_0,p_1,p_2$ are isomorphisms. Then $F$ is distinguished. Composing $S_X$ and $F$ makes $R$ distinguished; cancelling $S_Y$ makes $B$ distinguished. Applying Cart~(2) to $B$ and $p_1$ forces $p_3$ to be an isomorphism. This deduction concerns the abstract axioms. The editors' two counterexamples will show why the intended class cannot satisfy them.
\end{dcomment}


\noindent Ces axiomes sont dans la nature d'axiomes de stabilité, nous donnant des critères affirmant que certains carrés sont h-cartésiens, si telles conditions sont remplies (p. ex. que certains autres carrés sont déjà h-cartésiens). J'ai l'impression d'avoir épuisé ici les propriétés pertinentes de ce type (\textit{i.e.} que les autres devraient pouvoir s'en déduire). À présent, je m'occupe de propriétés d'existence et d'unicité de carrés h-cartésien. % [sic] carrés h-cartésien -> carrés h-cartésiens, missing plural agreement

\noindent\textbf{Cart (4) (Axiome du produit fibré).} Tout diagramme
\[
\begin{tikzcd}
X_1 \arrow[d] & \\
X_0 & X_2 \arrow[l]
\end{tikzcd}
\]
dans $\mathcal{H}$ peut se compléter en un carré h-cartésien
\[
\begin{tikzcd}
X_1 \arrow[d] & X_3 \arrow[l] \arrow[d] & \\
X_0 & X_2 \arrow[l] & .
\end{tikzcd}
\]
L'axiome suivant (\textit{cf.} son corollaire) assurera entre autres que le carré h-cartésien en question est \textit{unique à isomorphisme près} – mais on fera attention qu'il n'est en général pas unique “à isomorphisme unique près”.

% ---- from Der-ChXIV-07-10-10_p006-010.tex
\printedpage{4}
% head: A. Grothendieck : Dérivateurs, Ch. XIV, Édition des Universités de Montpellier II et Paris VII

\textbf{Cart (5) (Axiome d'exactitude)} Considérons deux carrés h-cartésiens
\[
X = \left(\begin{tikzcd}
X_1 \ar[d] & X_3 \ar[l] \ar[d] \\
X_0 & X_2 \ar[l]
\end{tikzcd}\right)
\qquad \text{et} \qquad
Y = \left(\begin{tikzcd}
Y_1 \ar[d] & Y_3 \ar[l] \ar[d] \\
Y_0 & Y_2 \ar[l]
\end{tikzcd}\right)
\]
et un homomorphisme $p = (p_0, p_1, p_2)$ entre les diagrammes de type
\[
\Phi \ = \
\begin{tikzcd}
1 \ar[d] & \\
0 & 2 \ar[l]
\end{tikzcd}
\]
associés (\textit{cf.} diagramme en traits pleins ci-dessous) :

\msp{page 5}
\[
\begin{tikzcd}[row sep=small, column sep=small]
& X_1 \ar[dl] \ar[dd, "p_1" near start] & & X_3 \ar[ll] \ar[dl] \ar[dd, dotted, "p_3"] \\
X_0 \ar[dd, "p_0"'] & & X_2 \ar[ll, crossing over] & \\
& Y_1 \ar[dl] & & Y_3 \ar[ll] \ar[dl] \\
Y_0 & & Y_2 \ar[ll] \ar[from=uu, crossing over, "p_2"' near end] & .
\end{tikzcd}
\]

Alors :
\begin{itemize}[leftmargin=2.5em]
\item[a)] Il existe un homomorphisme $\bar{p} : X \longrightarrow Y$ induisant $p$, \textit{i.e.} il existe une flèche $p_3$ de façon à compléter le diagramme cubique commutatif (en respectant donc la contrainte de commutativité pour les deux carrés adjacents à l'arête $p_3$).
\item[b)] Si de plus $p_0$, $p_1$, $p_2$ sont des isomorphismes, il en est de même de $p_3$ \texttt{[voir postface des éditeurs]}.
\end{itemize}

Prenant le cas où $p_0$, $p_1$, $p_2$ sont des identités, on trouve :

\textbf{Corollaire 1.} \textit{Sous les conditions de} Cart (4)\textit{, si on a} deux \textit{carrés h-cartésiens $X$, $X'$ qui complètent}
\[
\begin{tikzcd}[row sep=small]
& X_1 \ar[dl] & \\
X_0 & & , \\
& X_2 \ar[ul] &
\end{tikzcd}
\]
\textit{il existe un homomorphisme de $X$ dans $X'$ qui induise l'identité sur les sommets d'indices} 0, 1, 2\textit{, et un tel homomorphisme est toujours un isomorphisme.}

Mais cet isomorphisme n'est pas forcément unique, en d'autres termes, un automorphisme d'un carré h-cartésien peut induire l'identité sur les sommets $X_0$, $X_1$, $X_2$, sans être l'identité.
% foot: Version 7 octobre 2010

\dgrblock{dgrb744}{Existence of a comparison is not a chosen comparison}
\begin{dcomment}{exposition}{p5:sans être l'identité.}{Existence of a comparison is not a chosen comparison}
Cart~(4) supplies an object over each cospan. Cart~(5)(a) asks that a map of cospans extend to some map of these objects; it does not specify a choice compatible with composition. Cart~(5)(b), with its universal quantifier on the already chosen $p_3$, says that \emph{every} such extension over an isomorphism of cospans is invertible. Replacing ``every'' by ``there exists an invertible extension'' would be a different axiom.

For a fixed cospan $A\to C\leftarrow B$, consider the category of cones whose objects are triples $(T,a:T\to A,b:T\to B)$ satisfying $fa=gb$, and whose arrows preserve $a,b$. The first corollary says that the full subcategory consisting of distinguished completions is nonempty, that any two of its objects admit a comparison, and that every comparison is invertible. This is a connected groupoid, not a proof that it is a terminal category. Existence follows by Cart~(5)(a) with three identities; invertibility follows by (b). The source's nonuniqueness remark records the distinction between these conclusions. It supplies no separate example satisfying all five axioms with a nontrivial automorphism group.

For the intended homotopy class, the distinction is quite visible: the square with vertices $\Omega X,0,0,X$ and all edges zero permits every automorphism of $\Omega X$ over the fixed cospan. It also permits every \emph{noninvertible} endomorphism, including zero. Thus this example illustrates nonuniqueness but simultaneously invalidates (b); it cannot validate the literal axiom system. In an additive category satisfying the whole system, the development below proves that no such nonuniqueness is possible.

In the coherent repair the relevant base is itself a category of coherent cospans. Restriction from cartesian coherent squares to that category is an equivalence. This gives an honest functorial construction there, while forgetting coherence can again identify distinct maps. Neither a noncanonical choice of apex nor an arbitrary commuting cube restores the information that was forgotten.
\end{dcomment}


\printedpage{5}
% head: A. Grothendieck : Dérivateurs, Ch. XIV, Édition des Universités de Montpellier II et Paris VII

\msp{page 6}
\textbf{Corollaire 2.} \textit{Considérons un carré commutatif}
\[
Y = \left(\begin{tikzcd}
Y_1 \ar[d] & Y_3 \ar[l] \ar[d] \\
Y_0 & Y_2 \ar[l]
\end{tikzcd}\right) \quad ,
\]
\textit{et un objet $T$ de $\mathcal{H}$, d'où (posant} $\mathrm{Hom}(T,X) = X(T)$\textit{) un carré commutatif dans} (Ens)
\[
\begin{tikzcd}
Y_1(T) \ar[d] & Y_3(T) \ar[l] \ar[d] \\
Y_0(T) & Y_2(T) \ar[l]
\end{tikzcd} \quad ,
\]
\textit{et par suite}
\[
Y_3(T) \longrightarrow Y_1(T) \times_{Y_0(T)} Y_2(T) \quad . \tag*{$(*)$}
\]
\textit{Si $Y$ est h-cartésien, cette application} $(*)$ \textit{est} surjective \textit{quel que soit $T$.}

On applique Cart (5) à $Y$ et au carré constant $X$ de valeur $T$
\[
X = \quad
\begin{tikzcd}
T \ar[d, "\mathrm{id}"'] & T \ar[l, "\mathrm{id}"'] \ar[d, "\mathrm{id}"] \\
T & T \ar[l, "\mathrm{id}"']
\end{tikzcd} \quad ,
\]
en notant que la donnée de $p = (p_0, p_1, p_2)$ de
\[
\begin{tikzcd}
T \ar[d] & \\
T & T \ar[l]
\end{tikzcd}
\qquad\qquad \text{dans} \qquad\qquad
\begin{tikzcd}
Y_1 \ar[d] & \\
Y_0 & Y_2 \ar[l]
\end{tikzcd}
\]
revient à la donnée d'un élément du deuxième membre de $(*)$,
\[
\begin{tikzcd}
Y_1 \ar[d] & T \ar[l, "p_1"'] \ar[dl, "p_0"] \ar[d, "p_2"] \\
Y_0 & Y_2 \ar[l]
\end{tikzcd} \quad ,
\]
et que les éléments de $Y_3(T)$ qui s'envoient (par $(*)$) dans cet élément ne sont autres que les $p_3$ dont l'axiome Cart (5) affirme l'existence.

\textbf{N.B.} On notera que l'application surjective $(*)$ n'est en général \textit{pas} bijective, \textit{i.e.} le carré h-cartésien n'est en général pas cartésien.
% foot: Version 7 octobre 2010

\dgrblock{dgrb745}{Weak universality and the information in a homotopy}
\begin{dcomment}{exposition}{p6:le carré h-cartésien n'est en général pas cartésien.}{Weak universality and the information in a homotopy}
The second corollary identifies the positive content of the existence axiom. A distinguished square is a \emph{weak pullback}: every compatible pair of maps from $T$ factors through its apex, without uniqueness. The constant square on $T$ is distinguished by Cart~(2), and applying Cart~(5)(a) to it gives exactly the surjection displayed in the text. Conversely, if every distinguished square is a weak pullback, apply its weak property to $T=X_3$ and the two composites $X_3\to X_i\to Y_i$. This supplies Cart~(5)(a). Thus, given the constant squares, the Hom-set assertion and that part of Cart~(5) are equivalent.

For an actual homotopy pullback the surjection has a geometric explanation. In an enhanced homotopy theory, put
\[
 M_i=\operatorname{Map}(T,Y_i).
 \qquad
 \operatorname{Map}(T,Y_3)\simeq M_1\mathbin{\times^h_{M_0}}M_2.
\]
A point of the homotopy pullback consists of two maps and a path between their composites. Equality in $\pi_0M_0$ ensures that such a path exists, which proves surjectivity after taking components. Different choices of path can give different components. For the loop square, taking $T=S^0$ in pointed spaces and $X=S^1$ gives the map $\pi_0\Omega S^1\cong\mathbb Z\to\{*\}$, a fully explicit failure of injectivity. This statement is about the intended homotopy class, not an example of all five literal axioms. Mapping spaces supply the universal property before $\pi_0$ is taken; compare \citet[\S1.2.13]{Lurie2009HTT}.

There is a useful rigidity argument whenever an ordinary pullback $P=A\times_C B$ already exists in $\mathcal H$. Let $Q$ be a distinguished completion. Its canonical map $q:Q\to P$ has a section $s:P\to Q$ by weak universality. The endomorphism $sq$ of $Q$ is over the identity cospan, so Cart~(5)(b) makes it invertible. But $(sq)^2=sq$; an invertible idempotent is the identity. Hence $q$ is an isomorphism, and $Q$ is an ordinary pullback. Conversely, Cart~(4) and invariance under isomorphisms show that every existing ordinary pullback square is distinguished. Consequently, if $\mathcal H$ has all ordinary pullbacks, the proposed class is forced to be precisely that class. This is the familiar minimality argument for weak universal objects; its weak-kernel form is explicitly recorded by \citet[\S1, ``Weak kernels'', p.~6 of the 2014 manuscript]{Krause2016FlatCovers}.
\end{dcomment}


\printedpage{6}
% head: A. Grothendieck : Dérivateurs, Ch. XIV, Édition des Universités de Montpellier II et Paris VII

\msp{page 7}
\textbf{N.B.} Il y a une notion duale de \textit{catégorie à carrés h-cocartésiens,} ou \textit{catégorie quadrangulée à droite.} Le plus souvent, les deux structures sont présentes simultanément. Il y a lieu de dégager des propriétés qui relient les deux structures duales sur une même catégorie. Ainsi, sauf erreur, dans le cas $\mathcal{H}$ additive, les deux espèces de carrés, les h-cartésiens et les h-cocartésiens, devraient \textit{coïncider.} (Il me faudra revenir là-dessus par la suite.) Dans le cas $\mathcal{H}$ ``ponctuée'', il y a l'adjonction des foncteurs $\Sigma$, $\Omega$ liés l'un à la structure quadrangulée droite, l'autre à la structure quadrangulée gauche. Et il y a, quand il y a des \underline{Hom} internes ou externes, des propriétés éventuelles de $\uHom(X,Y)$ par rapport aux quadrangles h-cocartésiens en $X$, h-cartésiens en $Y$ -- mais c'est à peine un lien entre les deux structures \ldots

\dgrblock{dgrb746}{Additivity, stability, and the two variables of Hom}
\begin{dcomment}{later}{p7:mais c'est à peine un lien entre les deux structures}{Additivity, stability, and the two variables of Hom}
Grothendieck is looking for a common structure behind fibres and cofibres. Additivity alone cannot supply it. Take $\mathcal H=\mathrm{Vect}_k$, with ordinary pullback squares on the left and ordinary pushout squares on the right. Each structure satisfies the corresponding five axioms. For $V\ne0$, the square
\[
 \begin{tikzcd}[column sep=large]
 0\ar[d]&0\ar[l]\ar[d]\\
 V&0\ar[l]
 \end{tikzcd}
\]
is a pullback, since $0\times_V0=0$, but is not a pushout, since $0\amalg_00=0$. This refutes the additive speculation within the literal framework, without using any defective homotopy example.

For coherent diagrams, the corrected hypothesis is \emph{stability}. A pointed derivator is stable when its cartesian and cocartesian coherent squares coincide. The suspension--loop adjunction exists already in the pointed case; the equivalence of that adjunction characterizes stability \citetext{\citealp[Proposition~3.17]{Groth2013}; \citealp[Theorem~7.1, arXiv v2]{GrothPontoShulman2014}}. The corresponding higher-categorical statement appears in \citet[Proposition~1.1.3.4]{Lurie2017HA}. These theorems concern coherent squares, and do not rescue Cart~(3) on the ordinary homotopy category.

The Hom remark also has a precise scope. For fixed $T$, a derived mapping-space functor sends homotopy limits in its second variable to limits of spaces. In the first variable, mapping out of a homotopy colimit produces a homotopy limit. In a pointed enhancement this gives
\[
 \operatorname{Map}(\Sigma X,Y)
 \simeq\Omega\operatorname{Map}(X,Y)
 \simeq\operatorname{Map}(X,\Omega Y),
\]
where loops in the mapping space are based at the zero map; taking components gives the adjunction. Cisinski constructs the external derived Hom for model-derived derivators and proves its continuity \citep[Proposition~6.13, pp.~239--240]{Cisinski2003}. Internal Hom requires a specified closed monoidal structure and the appropriate derived adjunctions. Nothing on this page supplies such structure to an arbitrary category $\mathcal H$.
\end{dcomment}


\dgrblock{dgrb747}{Binary products force the literal structure to be ordinary}
\begin{dcomment}{development}{p7:mais c'est à peine un lien entre les deux structures}{Binary products force the literal structure to be ordinary}
The proposed axioms are meant to retain a weak universal property while allowing nonunique comparisons. The preceding rigidity argument shows that this freedom disappears wherever an ordinary pullback already exists. Pasting makes the restriction stronger: binary products alone force all the missing pullbacks to exist. This applies in particular to additive categories, without any assumption of kernels or of an abelian structure.

We assume Cart~(1)--(5) literally throughout this block. The weak-kernel minimality argument has the antecedent cited above. We checked the editors' postface, the coherent-square discussion of \citet[\S3]{Groth2013}, and \citet[\S1]{Krause2016FlatCovers}, together with targeted literature searches. We did not find the following deductions from these particular axioms there. The marking identifies proofs supplied by us, not a claim of priority.

\textbf{Pullback along a monomorphism.}\ours\quad
A distinguished square with a monomorphism among its two base arrows is an ordinary pullback.

\emph{Proof.} Write the square as the left square below, with $f$ monic. To test uniqueness of its lifting property, apply Cart~(4) to $P\xrightarrow qB\xleftarrow qP$ and obtain the right square:
\[
\begin{tikzcd}[column sep=large]
A\ar[d,"f"'] & P\ar[l,"p"']\ar[d,"q"]
 & R\ar[l,"r"']\ar[d,"s"]\\
C & B\ar[l,"g"] & P\ar[l,"q"] .
\end{tikzcd}
\]
Cart~(3) makes the outer rectangle distinguished. Its base cospan is $A\xrightarrow f C\xleftarrow{fp}P$, because $gq=fp$. Since $f$ is monic, this cospan has ordinary pullback $P$ with projections $p$ and $1_P$. The rigidity argument proved in the preceding exposition makes the canonical comparison $s:R\to P$ an isomorphism.

Weak universality of the right square supplies a diagonal $d:P\to R$ with $rd=sd=1_P$. Thus $d=s^{-1}$ and $r=s$. If $u,v:T\to P$ satisfy $qu=qv$, the same weak property gives $w:T\to R$ with $rw=u$ and $sw=v$; hence $u=v$. Therefore $q$ is monic. The original square is a weak pullback, and its factorization maps are unique because their composites with $q$ agree. It is an ordinary pullback. Cart~(1) treats the other base arrow.\hfill$\square$

\textbf{Collapse under binary products.}\ours\quad
If $\mathcal H$ has binary products and satisfies Cart~(1)--(5), it has all pullbacks, and its distinguished squares are exactly its ordinary pullback squares.

\emph{Proof.} Given $A\xrightarrow fC\xleftarrow gB$, form the cospan
\[
C\xrightarrow{\Delta}C\times C
  \xleftarrow{\,f\times g\,}A\times B.
\]
The diagonal is monic. A distinguished completion exists by Cart~(4), and the preceding result makes it an ordinary pullback. Its apex is $A\times_CB$: a pair $a:T\to A$, $b:T\to B$ with $fa=gb$ gives the unique map $(a,b)$ to $A\times B$ and the compatible map $fa$ to $C$, so the displayed pullback property gives a unique map to the apex. All pullbacks therefore exist. Rigidity shows that each distinguished square is ordinary; Cart~(4) and isomorphism invariance show that every ordinary pullback square is distinguished.\hfill$\square$

No terminal object was used. Additivity supplies the required binary products as biproducts, so the theorem is stronger than an additive statement. It does not classify categories lacking binary products, and we draw no general nonexistence conclusion in that setting.

\textbf{Pointed coincidence obstruction.}\ours\quad
If a pointed category carries both literal structures and their distinguished classes coincide, every object is isomorphic to zero.

\emph{Proof.} For any $X$, the cospan $0\to X\leftarrow0$ has ordinary pullback $0$. Cart~(4), rigidity, and isomorphism invariance designate the square with base $X$ and its other three vertices zero. Coincidence makes it distinguished cocartesian. Its upper span $0\leftarrow0\to0$ has ordinary pushout $0$; dual rigidity therefore gives $X\cong0$.\hfill$\square$

Even without coincidence, the literal loop construction is always zero, by the same pullback argument, and the literal suspension is zero if the right structure is present. Their adjunction would be the vacuous zero-functor adjunction. The meaningful suspension and loop functors in homotopy theory belong to the coherent setting anticipated on manuscript page~2.
\end{dcomment}


\section*{2\quad Exemples}

Un premier exemple trivial, est celui où $\mathcal{H}$ est stable par produits fibrés, et où les carrés h-cartésiens sont définis comme les carrés cartésiens. Les axiomes Cart (1) à Cart (5) sont alors des propriétés bien connues, $\pm$ tautologiques de la notion de carré cartésien. Et
\msp{page 8}
en un sens, le cas général doit se ramener à ce cas particulier banal.

Le seul type d'autre cas que je connaisse, est celui où $\mathcal{H} = \mathcal{M}W^{-1}$, où $W$ est un localiseur admettant une théorie de la fibre $W$-homotopique (\textit{cf.} chapitre XII). Dans ce cas, on l'a dit, on prend comme carrés h-cartésiens dans $\mathcal{H}$ ceux qui sont isomorphes, dans $\uHom(\bDelta_1 \times \bDelta_1, \mathcal{H})$, à un carré $W$-cartésien dans $\mathcal{M}$. Encore faut-il vérifier les axiomes Cart (1) à Cart (5). Je pense que je n'ai guère de chance d'y arriver que moyennant une

\textbf{Proposition 1.} \textit{Moyennant des hypothèses sur $(\mathcal{M}, W)$ qui vont être dégagées ci-dessous, vérifiées en tous cas si $(\mathcal{M}, W)$ est quillénisable, sur la catégorie des ensembles ordonnés finis} (\textit{cf.} XIII, p. 51)\textit{, on a ceci : soit}
\[
X = \left(\begin{tikzcd}
X_1 \ar[d] & X_3 \ar[l] \ar[d] \\
X_0 & X_2 \ar[l]
\end{tikzcd}\right)
\]
\textit{un carré dans $\mathcal{M}$, et soit} $\ho_W(X)$ \textit{son image dans} $\uHom(\bDelta_1 \times \bDelta_1, \mathcal{H})$\textit{. Pour que} $\ho_W(X)$ \textit{soit h-cartésien, il} faut (\textit{et il suffit, bien sûr}) \textit{que $X$ soit $W$-cartésien} \texttt{[voir postface des éditeurs]}\textit{.}
% foot: Version 7 octobre 2010

\dgrblock{dgrb748}{The precise assertion that fails in Proposition 1}
\begin{dcomment}{exposition}{p8:voir postface des éditeurs}{The precise assertion that fails in Proposition 1}
Grothendieck hopes to verify the axioms by transporting a homotopy-cartesian criterion through localization. The sufficient direction of Proposition~1 is the definition of $\operatorname{Cart}_{\mathrm{im}}$: a homotopy-cartesian strict square has a distinguished underlying image. The necessary direction would make homotopy cartesianness a property of the ordinary underlying square alone. The editors disprove precisely this reflection statement.

It matters that the obstruction occurs in Quillen model categories with entirely familiar diagram model structures. It cannot be removed by the promised availability of models for finite ordered indexing categories. The walking square and the walking cube are already finite posets; the example below lives in complexes over a field, where all their diagram model structures are available. An additional hypothesis could restrict the homotopy theory drastically, but Quillen presentability of diagrams does not establish the proposition.

The preliminary example, ordinary pullbacks in an ordinary category, remains valid. Symmetry, the isomorphism test, and pasting follow directly from the pullback universal property; every cospan has a completion by hypothesis; a map of cospans induces a unique map of pullbacks; an invertible base map induces an invertible comparison. These statements also prove the stronger uniqueness omitted from the proposed homotopical axioms. The preceding development gives a precise circumstance in which the general axiom system really does reduce to this example: the existence of binary products. It does not vindicate the intended model-category example.
\end{dcomment}


\printedpage{7}
% head: A. Grothendieck : Dérivateurs, Ch. XIV, Édition des Universités de Montpellier II et Paris VII

\section*{Postface des éditeurs}

La proposition ci-dessus est fausse. On rappelle que si $\mathcal{M}$ est une catégorie de modèles fermée de Quillen de classe d'équivalences faibles $W$, un carré commutatif
\[
X = \left(\begin{tikzcd}
X_1 \ar[d] & X_3 \ar[l] \ar[d] \\
X_0 & X_2 \ar[l]
\end{tikzcd}\right)
\]
de $\mathcal{M}$ est $W$-cartésien si et seulement si il existe un carré \textit{cartésien}
\[
Y = \left(\begin{tikzcd}
Y_1 \ar[d] & Y_3 \ar[l] \ar[d] \\
Y_0 & Y_2 \ar[l]
\end{tikzcd}\right)
\]
avec $Y_1 \longrightarrow Y_0$ une fibration, et $Y_0$ et $Y_2$ fibrants (cette dernière condition étant superflue si la catégorie de modèles $\mathcal{M}$ est propre à droite), et un morphisme de carrés $s : X \longrightarrow Y$, $s = (s_0, s_1, s_2, s_3)$, \textit{i.e.} un cube commutatif
\[
\begin{tikzcd}[row sep=small, column sep=small]
& X_1 \ar[dl] \ar[dd, "s_1" near start] & & X_3 \ar[ll] \ar[dl] \ar[dd, "s_3"] \\
X_0 \ar[dd, "s_0"'] & & X_2 \ar[ll, crossing over] & \\
& Y_1 \ar[dl] & & Y_3 \ar[ll] \ar[dl] \\
Y_0 & & Y_2 \ar[ll] \ar[from=uu, crossing over, "s_2"' near end] & ,
\end{tikzcd}
\]
tel que $s$ soit une équivalence faible argument par argument, \textit{i.e.} $s_0, s_1, s_2, s_3 \in W$. Il est équivalent de demander que pour \textit{tout} diagramme commutatif à flèches verticales des équivalences faibles
\[
\begin{tikzcd}
X_1 \ar[r] \ar[d] & X_0 \ar[d] & X_2 \ar[l] \ar[d] \\
Y_1 \ar[r] & Y_0 & Y_2 \ar[l]
\end{tikzcd}
\]
tel que $Y_1 \longrightarrow Y_0$ soit une fibration, et $Y_0$, $Y_2$ des objets fibrants, la flèche canonique $X_3 \longrightarrow Y_1 \times_{Y_0} Y_2$ soit une équivalence faible.

\dgrblock{dgrb749}{Postface: why one strict comparison cube suffices}
\begin{dcomment}{exposition}{p8:Il est équivalent de demander}{Postface: why one strict comparison cube suffices}
\emph{Editors' postface, printed p.~7.}
The editors replace the cospan before taking its limit. This makes the missing comparison map canonical, and explains the difference between a weak equivalence of whole squares and a mere isomorphism after applying $\operatorname{dia}_{\square}$.

Here is a verification with exactly their fibrancy hypotheses. Give $\Phi=(1\to0\leftarrow2)$ the Reedy degrees $\deg(2)=0$, $\deg(0)=1$, $\deg(1)=2$, with $2\to0$ direct and $1\to0$ inverse. A diagram $Y$ is Reedy fibrant precisely when $Y_0,Y_2$ are fibrant and $Y_1\to Y_0$ is a fibration. The constant-diagram functor is left Quillen: for a (trivial) cofibration its relative latching maps at $1,2$ are that map, and at $0$ the identity. Hence its right adjoint, the pullback functor, is right Quillen. This is the mixed Reedy method of \citet[Theorem~5.2.5 and the proof of Lemma~5.2.6]{Hovey1999ModelCategories}, dualized.

For a cospan $K$, choose a Reedy trivial cofibration $K\to RK$ with fibrant target. Any objectwise weak equivalence $K\to Y$ to a cospan with the printed fibrancy conditions extends to a map $RK\to Y$ by lifting against $Y\to *$. It is a weak equivalence by two-out-of-three. Applying the right Quillen pullback functor preserves it, since both diagrams are fibrant. Consequently the two comparison maps from the original fourth vertex are weak equivalences simultaneously. This proves ``some'' if and only if ``every''. Taking $Y_3=\lim Y$ makes the fourth component of the strict cube the unique induced cone map, so it also proves the single-cube criterion.

For completeness, right properness removes both objectwise fibrancy conditions. If $K_1\to K_0$ is a fibration and $K\to R$ is a fibrant replacement of the cospan, put $K'_1=R_1\times_{R_0}K_0$. Right properness and two-out-of-three make $K_1\to K'_1$ a weak equivalence between fibrant objects of the slice $\mathcal M/K_0$. Pullback to $K_2$ is right Quillen, so
\[
K_1\times_{K_0}K_2\longrightarrow
R_1\times_{R_0}K_2\longrightarrow R_1\times_{R_0}R_2
\]
are weak equivalences: the first by that slice argument, the second by right properness. This verifies the parenthesis in the source, without adding a properness assumption to the fibrant version.
\end{dcomment}


Supposons maintenant que la catégorie de modèles $\mathcal{M}$ soit ponctuée, autrement dit qu'elle admette un objet $*$ à la fois initial et final, et soit $X$ un objet fibrant de $\mathcal{M}$. Choisissons une décomposition de la flèche $* \longrightarrow X$, en une équivalence faible suivie d'une fibration
\[
* \longrightarrow C \longrightarrow X \quad ,
\]
et formons le carré cartésien
% foot: Version 7 octobre 2010
\printedpage{8}
% head: A. Grothendieck : Dérivateurs, Ch. XIV, Édition des Universités de Montpellier II et Paris VII
\[
\begin{tikzcd}
C \ar[d] & \Omega X \ar[l] \ar[d] \\
X & * \ar[l]
\end{tikzcd} \quad ,
\]
où par définition $\Omega X = C \times_X *$ est l'\textit{espace des lacets} de $X$. Ce carré est $W$-cartésien. Par ailleurs, il est isomorphe dans $\uHom(\bDelta_1 \times \bDelta_1,\ \mathcal{M}W^{-1})$ au carré trivial
\[
\begin{tikzcd}
* \ar[d] & \Omega X \ar[l] \ar[d] \\
X & * \ar[l]
\end{tikzcd} \quad .
\]

Néanmoins, ce dernier n'est \textit{pas} en général $W$-cartésien. En effet, s'il l'était, l'unique flèche $s_3 : \Omega X \longrightarrow \Omega X$ de $\mathcal{M}$ rendant commutatif le cube
\[
\begin{tikzcd}[row sep=small, column sep=small]
& * \ar[dl] \ar[dd] & & \Omega X \ar[ll] \ar[dl] \ar[dd, dotted, "s_3"] \\
X \ar[dd, "\parallel"'] & & * \ar[ll, crossing over] & \\ % Audit resolved: vertical identity mark on X; see corrections.csv, PDF leaf 10.
& C \ar[dl] & & \Omega X \ar[ll] \ar[dl] \\
X & & * \ar[ll] \ar[from=uu, crossing over] &
\end{tikzcd}
\]
serait une équivalence faible. Or, comme la flèche $s_3$ se factorise par $* \times_X * \simeq *$, cela impliquerait que $\Omega X \to *$ est une équivalence faible. Cette condition est très restrictive. Si, par exemple, $\mathcal{M}$ est la catégorie de modèles des espaces topologiques pointés, et $X$ un CW-complexe pointé, elle équivaut à la contractibilité de la composante connexe de $X$ contenant le point base. De même, si $\mathcal{M}$ est la catégorie de modèles des complexes d'une catégorie abélienne de Grothendieck et $X$ un complexe, elle signifie que $X$ est acyclique.

\dgrblock{dgrb750}{Postface: a two-term complex and the missing filling}
\begin{dcomment}{examples}{p8:Cette condition est très restrictive.}{Postface: a two-term complex and the missing filling}
\emph{Editors' postface, printed p.~8.}
The counterexample is especially transparent when every map can be written as a scalar. Let $k$ be a nonzero field and use unbounded cochain complexes, with quasi-isomorphisms as weak equivalences and degreewise surjections as fibrations \citep[Theorem~2.3.11, with degrees reversed]{Hovey1999ModelCategories}. Our shift convention is $(V[r])^n=V^{n+r}$, with differential $(-1)^r d_V$.

Take $X=k[0]$ and the acyclic complex $C=(k\xrightarrow{1}k)$ in degrees $0,1$. Let $p:C\to X$ be the identity in degree $0$ and zero in degree $1$. This is a surjective chain map. Its kernel $P$ is $k[-1]$, concentrated in degree $1$. Thus $0\xrightarrow{\sim}C\xrightarrow pX$ is the editors' factorization and
\[
\begin{tikzcd}[column sep=large]
C\ar[d,"p"'] & P\ar[l,"i"']\ar[d]\\
X & 0\ar[l]
\end{tikzcd}
\]
is both an ordinary and a homotopy pullback. The contraction $s^1:C^1\to C^0$ is the identity, and $d_Cs+sd_C=1_C$. In $D(k)$, therefore, $C\cong0$, and the underlying square is isomorphic to the square with apex $P$, base $X$, and all edges zero.

That latter square, taken \emph{strictly} in complexes, has comparison $s_3=0_P$ to the displayed fibrational replacement: commutativity requires $is_3=0$, and $i$ is injective. The comparison induces zero on $H^1(P)=k$, so it is not a quasi-isomorphism. This computes the dotted arrow in the postface and disproves the necessary direction of Proposition~1.

To locate the information lost, write down a homotopy-pullback model for $f:A\to Z\leftarrow B:g$:
\[
\begin{split}
H(f,g)^n&=A^n\oplus B^n\oplus Z^{n-1},\\
d(a,b,z)&=(d_Aa,d_Bb,fa-gb-d_Zz).
\end{split}
\]
The chain-map identities imply $d^2=0$. This is a shifted mapping cone of $f-g$, with the last coordinate negated relative to one common convention. A cone with maps $u:T\to A$, $v:T\to B$ and a specified degree $-1$ homotopy $h$ satisfying
\[
 d_Zh+hd_T=fu-gv
\]
has comparison $c_h=(u,v,h):T\to H(f,g)$. The displayed identity is exactly its chain-map condition. The coherent square is cartesian precisely when this comparison is a quasi-isomorphism.

For $A=B=0$, $Z=X=k[0]$, and $T=P=k[-1]$, the four edges are zero and $H(0,0)=P$. Every scalar $\lambda\in k$ specifies a filling $h_\lambda^1=\lambda$, and $c_{h_\lambda}$ is multiplication by $\lambda$ on $P$. It is a quasi-isomorphism exactly when $\lambda\ne0$. Its mapping cone is the two-term complex
\[
 k\xrightarrow{\lambda}k
 \quad\hbox{in degrees }0,1,
\]
which is acyclic for $\lambda\ne0$ and has $k$ in both cohomology degrees for $\lambda=0$. Thus identical vertices and edges support both a cartesian and a noncartesian coherent filling.

We can check the transport sign as well. Put $t=ps$, so $d_Xt+td_C=p$. The map
\[
 H(p,0)\longrightarrow X[-1],\qquad (c,z)\longmapsto z-tc
\]
is a chain map with acyclic kernel, the diagonal copy of $C$. Composing with $P\to H(p,0)$, $a\mapsto(i(a),0)$, gives $-1_P$. The strict pullback transported through this contraction has filling $h_{-1}$, whereas the strict zero-edge square has filling $h_0$. The sign depends on the displayed cone convention; its being nonzero does not. With the boundary and its identifications fixed, the classes of fillings are indexed by $k$, because $\underline{\operatorname{Hom}}(P,X)$ is concentrated in degree $-1$ and has no degree $-2$ boundaries. Allowing scalar automorphisms of the apex identifies all nonzero fillings, while the zero filling remains separate.

The editors' other interpretations agree with this computation. In pointed spaces, $X=S^1$ gives $\pi_0\Omega X=\mathbb Z$, so the zero self-map cannot be a weak equivalence. For a CW complex in general, the homotopy groups of its loop space detect all positive homotopy groups of the base component; Whitehead's theorem then makes exactly that component contractible \citep[Theorem~4.5 and \S4.3]{Hatcher2002}. Other components are unrestricted. For unbounded complexes in a Grothendieck abelian category, an acyclic middle term in $0\to P\to C\to X\to0$ gives $H^n(P)\cong H^{n-1}(X)$, so $P$ is acyclic exactly when $X$ is. The standard injective model structure supplies this setting \citep[Theorem~2.2 and Proposition~2.12 of the preprint]{Hovey2001Sheaves}; no assertion that acyclic complexes are chain-contractible is needed.
\end{dcomment}


Des considérations analogues montrent que les carrés $h$-cartésiens de $\mathcal{M}W^{-1}$, définis dans la section 2 comme étant les carrés commutatifs de $\mathcal{M}W^{-1}$ qui sont isomorphes dans $\uHom(\bDelta_1 \times \bDelta_1,\ \mathcal{M}W^{-1})$ à un carré $W$-cartésien, ne satisfont pas à toutes les conditions de Cart (3) et Cart (5) de la section 1. En effet, supposons toujours que $\mathcal{M}$ soit ponctuée, et gardons les notations ci-dessus. Considérons le diagramme commutatif dans $\mathcal{M}W^{-1}$
\[
\begin{tikzcd}
* \ar[d] & \Omega X \ar[l] \ar[d] & \Omega X \ar[l, "*_{\Omega X}"'] \ar[d] \\
X & * \ar[l] & * \ar[l]
\end{tikzcd} \quad ,
\]
où $*_{\Omega X}$ désigne l'unique endomorphisme de $\Omega X$ se factorisant par $*$. Il résulte de ce qui précède que le carré de gauche et le carré composé sont $h$-cartésiens. En revanche, le carré de droite ne l'est pas en général. En effet, si ce carré était isomorphe dans $\uHom(\bDelta_1 \times \bDelta_1,\ \mathcal{M}W^{-1})$ à un carré $W$-cartésien, la flèche $*_{\Omega X}$ serait un isomorphisme
% foot: Version 7 octobre 2010

% ---- from Der-ChXIV-07-10-10_p011-011.tex
\printedpage{9}
% head: A. Grothendieck : Dérivateurs, Ch. XIV, Édition des Universités de Montpellier II et Paris VII
% foot: Version 7 octobre 2010
de $\cM W^{-1}$, donc une équivalence faible de $\cM$, ce qui impliquerait les mêmes conditions très restrictives que ci-dessus. Ainsi l'une des implications de l'axiome Cart 3 n'est pas satisfaite.

De même le cube commutatif dans $\cM W^{-1}$
\[
\begin{tikzcd}[row sep=small, column sep=small]
 & * \arrow[dl] \arrow[dd] & & \Omega X \arrow[ll] \arrow[dl] \arrow[dd, "*_{\Omega X}"] \\
X \arrow[dd, "\shortparallel"'] & & * \arrow[ll, crossing over] & \\
 & * \arrow[dl] & & \Omega X \arrow[ll] \arrow[dl] \\
X & & * \arrow[ll] \arrow[from=uu, crossing over] & \mathllap{,}
\end{tikzcd}
\]
dont les faces horizontales sont $h$-cartésiennes, montre que la condition $(b)$ de l'axiome Cart (5) n'est pas satisfaite en général.


\dgrblock{dgrb751}{Postface: the two failed quantifiers}
\begin{dcomment}{status}{p8:dont les faces horizontales sont}{Postface: the two failed quantifiers}
\emph{Editors' postface, printed pp.~8--9.}
The two final counterexamples use the same object $P=\Omega X$ and the same underlying zero-edge square $Q$. They apply whenever $P$ is not isomorphic to zero, a condition realized by both worked examples above. In the horizontal pair, the left square and outer rectangle are $Q$, hence belong to $\operatorname{Cart}_{\mathrm{im}}$. In the right square the bottom edge is $1_0$ and the top edge is $0_P$. A homotopy pullback of an equivalence has an equivalence as opposite edge, so this right square cannot belong to $\operatorname{Cart}_{\mathrm{im}}$. This is exactly the implication ``I and III cartesian imply II cartesian'' that fails.

For Cart~(5)(b), use the endomorphism of $Q$ that is the identity at the three cospan vertices and zero at $P$. It is a commuting cube in the ordinary homotopy category, but its fourth component is not invertible. The positive existence statement in (a), and the failure of uniqueness already noticed in Corollary~2, should not be conflated with this false assertion about \emph{every} fourth component. As the formal deduction above shows, (b) would have followed from the full pasting axiom anyway.

The implication $0_P\in W\Rightarrow P\to0\in W$ used in the postface needs no cofibrancy assumption. Writing $r:P\to0$ and $s:0\to P$, the arrow $r$ is a retract, in the arrow category, of $sr=0_P$, using components $(1_P,s)$ and $(1_P,r)$. Weak equivalences are closed under retracts. The converse follows from two-out-of-three. Equivalently, an invertible zero endomorphism makes the object zero in the homotopy category. These are settled counterexamples, not conjectures awaiting a stronger model-category theorem.
\end{dcomment}


\dgrblock{dgrb752}{Why composition nevertheless survives}
\begin{dcomment}{exposition}{p8:dont les faces horizontales sont}{Why composition nevertheless survives}
The failure of cancellation leaves a useful positive statement: for every Quillen model category, $\operatorname{Cart}_{\mathrm{im}}$ is closed under composition. Two adjacent cartesian witnesses can be made compatible along their shared edge. A left square and an outer rectangle, as in the counterexample, are different data. We give the rectification argument to keep this distinction precise.

For any model category $\mathcal A$, the arrow comparison
\[
U_{\mathcal A}:\operatorname{Ho}(\mathcal A^{[1]})
 \longrightarrow\operatorname{Ho}(\mathcal A)^{[1]}
\]
is full, essentially surjective, and conservative. The first two properties are the arrow-shape lifting property used in derivator strength \citep[Lemma~1.9]{Groth2013}. To verify fullness directly, take a Reedy-cofibrant source arrow $f:X_0\to X_1$ and a Reedy-fibrant target arrow $g:Y_0\to Y_1$. Thus $X_0$ is cofibrant, $f$ is a cofibration, and $Y_0,Y_1$ are fibrant. Represent a square in the homotopy category by maps $a_0,a_1$ and a right homotopy between $a_1f$ and $ga_0$. The endpoint map $\mathrm{ev}_0:\operatorname{Path}(Y_1)\to Y_1$ is a trivial fibration. Lifting against $f$ extends that homotopy from $X_0$ to $X_1$, replacing $a_1$ by a homotopic $a'_1$ with $a'_1f=ga_0$ strictly. This gives the required coherent-arrow morphism. Any arrow in the homotopy category has an actual representative between cofibrant-fibrant objects, proving essential surjectivity. Finally, pointwise weak equivalences and saturation imply conservativity. In particular, an isomorphism of underlying arrows lifts to an isomorphism of coherent arrows.

Now choose cartesian strict lifts of the two adjacent squares. Their common columns may differ, but the given underlying identifications specify an isomorphism between these columns in $\operatorname{Ho}(\mathcal M)^{[1]}$. Apply the preceding result for $\mathcal A=\mathcal M$ to lift it. The two squares can then be regarded as composable morphisms in $\operatorname{Ho}(\mathcal M^{[1]})$. Choose bifibrant representatives for the three column objects in the arrow model category, with a single shared middle representative, and represent the two morphisms by actual maps. They form one strict rectangle in $\mathcal M$.

Apply the arrow comparison a second time, now to $\mathcal A=\mathcal M^{[1]}$. Each new small square has an underlying arrow in $\operatorname{Ho}(\mathcal A)$ isomorphic to that of its original witness. Fullness lifts this isomorphism and conservativity makes the lift invertible in $\operatorname{Ho}(\mathcal M^{\square})$. Thus the two new squares are cartesian as coherent squares. Coherent pasting makes their outer rectangle cartesian, and its underlying square is isomorphic to the required composite. This proves the assertion without a combinatorial or stability assumption, and without lifting an arbitrary incoherent cube.

In the stable setting, \citet[Proposition~6.11, pp.~17--18 of the 7~June~2021 author version]{ChristensenFrankland2022} prove the corresponding pasting statement even for homotopy-cartesian squares in an arbitrary triangulated category. The valid remnant of the proposed calculus is therefore substantial: symmetry, the isomorphism test, completion of cospans, weak universality, and positive composition survive for the homotopy-image class. Cancellation and Cart~(5)(b) do not.
\end{dcomment}


\dgrblock{dgrb753}{The coherent replacement of the square axioms}
\begin{dcomment}{later}{p8:dont les faces horizontales sont}{The coherent replacement of the square axioms}
\emph{After the editors' postface.}
The construction Grothendieck names on manuscript page~2 gives the appropriate replacement. Let $\mathbb D$ be a derivator and $j:\Phi\hookrightarrow\square$ the full inclusion omitting the prospective pullback vertex. A coherent square $Q\in\mathbb D(\square)$ is cartesian when the adjunction unit
\[
 Q\longrightarrow j_*j^*Q
\]
is invertible. Right Kan extension along $j$ is fully faithful; consequently
\[
 j^*:\mathbb D(\square)_{\mathrm{cart}}
       \simeq\mathbb D(\Phi)
\]
is an equivalence, with inverse $j_*$ \citep[Proposition~1.20 and Definition~3.9]{Groth2013}. A coherent cospan therefore has a cartesian completion, and a morphism of coherent cospans has a unique extension between these completions once their identifications are fixed. If its three components are isomorphisms, pointwise conservativity makes it an isomorphism in $\mathbb D(\Phi)$; the equivalence makes the extended square morphism invertible. This proves the correct fourth-vertex assertion. Its input is a morphism in the coherent category, not an arbitrary cube in $\mathbb D(e)$.

The full pasting-and-cancellation law is a statement about one coherent rectangle. To see the mechanism in a model or higher-category enhancement, write it as
\[
\begin{tikzcd}
A\ar[d]&B\ar[l]\ar[d]&C\ar[l]\ar[d]\\
D&E\ar[l]&F\ar[l].
\end{tikzcd}
\]
If the left square is cartesian, $B\simeq A\times^h_D E$. Associativity of the homotopy limit gives
\[
 B\times^h_E F\simeq
 (A\times^h_D E)\times^h_E F\simeq A\times^h_D F.
\]
Coherence identifies the comparison from $C$ to the first expression with its comparison to the last. One is an equivalence exactly when the other is. For derivators this is the right-Kan-extension pasting statement of \citet[Proposition~3.13]{Groth2013}; the higher-categorical proof is given, dually, by \citet[Lemma~4.4.2.1]{Lurie2009HTT}. The editors' rectangle loses the compatibility needed to make those two comparison maps agree.

In a stable enhancement there is a useful translation into the underlying triangulated category. A coherent cartesian square with projections $u:P\to A$, $v:P\to B$ and base maps $f:A\to C$, $g:B\to C$ supplies a distinguished triangle
\[
 P\xrightarrow{(u,-v)}A\oplus B
 \xrightarrow{(f,g)}C\xrightarrow{\delta}\Sigma P
\]
\citep[Theorem~6.1, arXiv v2]{GrothPontoShulman2014}. The connecting morphism is part of this conclusion. In the zero-edge loop example it is an isomorphism $X\to\Sigma P$, with the sign dictated by conventions. The cube with fourth component $0_P$ and base component $1_X$ cannot preserve that morphism, since it would require $\delta=(\Sigma0_P)\delta=0$. The triangle therefore diagnoses exactly why that cube is incoherent, even though its four visible faces commute in the homotopy category.

The relation with higher categories is a comparison of enhanced diagram theories. For a suitable model presentation, localizing diagram categories gives the homotopy categories of coherent diagrams; precise comparisons appear in \citet[Propositions~1.3.4.24--1.3.4.25]{Lurie2017HA}. It is not necessary here to identify every abstract derivator with an infinity category. The correct conclusion is narrower and sufficient: the familiar homotopy theories retain the diagram information on which the failed inference depended.
\end{dcomment}
\begin{otherlanguage}{english}
\dgroverview{XIV}{Chapter XIV \textperiodcentered\ \textup{XIV} \textperiodcentered\ Overview}
\selectlanguage{english}
Chapter XIV is a short experiment in reducing homotopy theory to a calculus of squares in a single category. Its eight manuscript pages introduce five proposed axioms and begin the examples. The editors' postface, which has no manuscript pagination, explains why the intended model-category example fails. The fragment is especially instructive because Grothendieck identifies the necessary finer setting before attempting the reduction: square diagrams must be localized as diagrams, together with their maps, rather than reconstructed from four arrows in the localized category.

The mathematical object under discussion can be described without ambiguity. For $\mathcal H=\mathcal MW^{-1}$, let $\operatorname{Cart}_{\mathrm{im}}$ consist of those commutative squares in $\mathcal H$ that are isomorphic to the underlying images of homotopy-cartesian strict squares in $\mathcal M$. With $\square$ denoting the walking commutative square, the relevant functor is
\[
\operatorname{dia}_{\square}:
\mathcal M^{\square}(W^{\square})^{-1}
\longrightarrow\mathcal H^{\square}.
\]
The domain retains coherent diagrams. The target retains ordinary commutativity after localization. Membership in $\operatorname{Cart}_{\mathrm{im}}$ is an existential statement about a coherent realization. Proposition~1 would promote it to a statement about every strict realization, and this is the false inference.

There is a specific historical reason for wanting the larger diagram theory. In \emph{Pursuing Stacks}, \S69, Grothendieck connects it with the trace and determinant problems that motivated him in the period of SGA~5; the letter to Thomason of 2~April~1991 again stresses derived diagram categories and mapping cones \citep{Grothendieck1983,Grothendieck1991Letter}. The present chapter tests whether one can retain a useful part of their calculus in a simpler language. The editors' preliminary notice suggests that recognition of the failed example explains the abrupt ending. That suggestion is conjectural; the postface establishes the mathematical failure, not the author's reason for interrupting the manuscript.

The formal axioms are consistent: any category with pullbacks, equipped with its ordinary pullback squares, satisfies them. Cart~(1) gives symmetry, Cart~(2) tests trivial squares against isomorphisms, and Cart~(3) combines pasting with cancellation. Cart~(4) requires a completion of every cospan. Cart~(5)(a) extends a map between the three-vertex boundaries of two distinguished squares; (b) asserts that every such extension is invertible whenever the boundary map is invertible. The commentary supplies the cube proof, including the stronger three-face hypothesis in the footnote, and proves that (b) already follows from Cart~(1)--(3).

The source's second corollary isolates the weak universal property that does survive in homotopy theory. For a homotopy pullback $P$ of $A\to C\leftarrow B$, there is a surjection
\[
[T,P]\longrightarrow[T,A]\times_{[T,C]}[T,B].
\]
A compatible pair in the target admits a homotopy between its composites; a lift remembers a choice of that homotopy. Passing to components forgets the choice, so the map need not be injective. In the loop square for $S^1$, with $T=S^0$ in pointed spaces, it is the map $\mathbb Z\to\{*\}$. The surjection is equivalent to Cart~(5)(a) for a class containing the constant squares. The much stronger assertion in (b) does not follow from this weak property alone.

Positive composition does hold for the homotopy-image class of every Quillen model category. The commentary proves this by applying fullness and conservativity of the arrow comparison twice, first to the model category and then to its arrow category, so that two cartesian witnesses acquire a single shared middle column. This is compatible with the postface: its counterexample concerns cancellation.

The editors' example distinguishes the underlying square from its realization. Factor the basepoint of a fibrant object as
\[
0\xrightarrow{\sim}C\twoheadrightarrow X,
\qquad P=C\times_X0.
\]
The resulting strict pullback is homotopy cartesian. Its image in $\mathcal H^{\square}$ is isomorphic to the square with apex $P$, base $X$, the other vertices zero, and all edges zero. But that zero-edge square, regarded strictly in $\mathcal M$, has comparison $0_P:P\to P$ to the fibrational replacement. It is homotopy cartesian exactly when $P$ is weakly zero. This proves the failure of Proposition~1 whenever the loop object is nontrivial. The same underlying square gives the two final counterexamples: a horizontal rectangle with left and outer squares distinguished but right square not distinguished, and a cube whose first three components are identities while the fourth is zero.

The worked algebraic example takes $X=k[0]$ and $C=(k\xrightarrow1k)$ in degrees $0,1$. Then $P=k[-1]$. A specified filling of the fixed zero-edge square is a scalar $\lambda$; the homotopy-pullback comparison is multiplication by $\lambda$ on $P$. Its cone is $k\xrightarrow\lambda k$, so the coherent square is cartesian precisely for $\lambda\ne0$. The commentary computes the transport through a contraction and its sign explicitly. It also checks the editors' topological qualification: for a CW complex, weak triviality of the loop object means contractibility of the base component, with no restriction on other components. For unbounded complexes in a Grothendieck abelian category it means acyclicity.

The literal axioms have a stronger limitation than the displayed counterexample alone reveals. Our development proves that, if $\mathcal H$ has binary products, the axioms force all ordinary pullbacks to exist and force the distinguished class to be exactly their class. The proof first compares a distinguished weak pullback with an existing ordinary one: a section and the invertibility of a cone-preserving idempotent make the comparison an isomorphism. This is the minimality argument familiar for weak kernels \citep[\S1, p.~6 of the 2014 manuscript]{Krause2016FlatCovers}. Pasting a distinguished square with a weak kernel pair then shows that pullbacks along monomorphisms are ordinary. Finally one pulls back the diagonal $C\to C\times C$ along $A\times B\to C\times C$. The complete argument is given in the marked development block; its application to these axioms was not located in the bounded literature search, and no priority is claimed.

Additive categories satisfy the binary-products hypothesis. Moreover, even without products, the literal loop object in a pointed category must be zero, because $0\to X\leftarrow0$ already has ordinary pullback zero. If the dual structure also exists, its suspension is zero. If the two classes coincide, the square with base $X$ and three zero vertices makes $X$ itself zero. Thus Grothendieck's additive coincidence speculation is false for the literal structures: ordinary pullbacks and pushouts in vector spaces give an immediate counterexample.

The modern positive statement belongs to coherent squares. For the inclusion $j:\Phi\hookrightarrow\square$ of the cospan, cartesian coherent squares are the essential image of the fully faithful right Kan extension $j_*$. Restriction consequently identifies their category with $\mathbb D(\Phi)$. This makes the completion functorial and proves the correct fourth-vertex assertion for coherent morphisms. Coherent rectangles admit pasting and cancellation. In a pointed derivator the suspension--loop adjunction exists; stability, rather than additivity, is the condition under which cartesian and cocartesian coherent squares agree, equivalently that adjunction is an equivalence \citetext{\citealp[Proposition~3.17]{Groth2013}; \citealp[Theorem~7.1, arXiv v2]{GrothPontoShulman2014}}. External derived Hom has the expected continuity, with the usual enrichment hypotheses for an internal version \citep[Proposition~6.13]{Cisinski2003}. None of these assertions makes the forgetful functor reflect the property disproved in the postface.

The proposed main example and Proposition~1 are therefore settled negatively, while the coherent replacements are established theorems. The manuscript's broad request for further compatibility axioms does not by itself specify a present open problem. This commentary does not classify all literal quadrangulated categories without binary products or produce a nonordinary example in that generality. Its audit keeps that limit separate from the explicit counterexamples and the proved structural deductions.
\end{otherlanguage}
\dgrchunkend
